
\chapter{Первый автомат: Lamp}
\label{chap:first-automaton}

Начнём с простой лампы, которую можно включать и выключать.
У неё два состояния: \code{off} и \code{on}. Команды
\code{switch_on} и \code{switch_off} переводят лампу
из одного состояния в другое.

Дополнительно будем считать, сколько раз лампу включали.

\begin{figure}[ht]
  \centering
  \begin{tikzpicture}[
  state/.style={draw=ciaoBlue,fill=ciaoPaleBlue!45,very thick,
    rounded corners=2mm,minimum width=32mm,minimum height=15mm,
    font=\sffamily\bfseries},
  flow/.style={-{Stealth[length=3mm]},thick,draw=ciaoNavy},
  every node/.style={font=\small}
]
  \node[state] (off) {off};
  \node[state,right=55mm of off] (on) {on};
  \draw[flow] (off) to[bend left=18]
    node[above,fill=white,inner sep=2pt]{switch\_on} (on);
  \draw[flow] (on) to[bend left=18]
    node[below,fill=white,inner sep=2pt]{switch\_off} (off);
  \draw[flow] ($(off.west)+(-14mm,0)$) -- (off.west);
\end{tikzpicture}

  \caption{Состояния и переходы лампы}
  \label{fig:lamp}
\end{figure}

\section{Описание автомата}

Автомат объявляется как Python-класс, наследующийся
от \code{Automaton}. Внутри него определяются состояния,
команды, действия и переходы между состояниями.

\lstinputlisting[
  caption={Объявление автомата Lamp},
  label={lst:lamp}
]{code/lamp.py}

\subsection{Состояния и контекст}

Состояния \code{off} и \code{on} определяют,
включена лампа или выключена.
Параметр \code{initial=True} указывает начальное
состояние автомата.

Помимо состояний, автомат может хранить дополнительные
данные в контексте. В нашем случае \code{LampContext}
содержит поле \code{switch_on_count} --- количество
включений лампы.

Таким образом, состояние определяет текущее поведение
автомата, а контекст хранит связанные с ним данные.

\subsection{Команды и действия}

Команды \code{switch_on} и \code{switch_off} объявлены
с помощью \code{@provided_command}. Они определяют
события, на которые может реагировать автомат.

Тела команд оставлены пустыми намеренно: команды
не выполняют действия напрямую, а запускают переходы.
Параметр \code{external=True} позволяет отправлять
их извне программы.

Для изменения данных контекста используются действия,
отмеченные декоратором \code{@action}.

В нашем примере \code{increment_switch_on_count}
увеличивает счётчик включений. Параметр \code{ctx}
предоставляет доступ к контексту автоматного объекта.

Необязательный аргумент \code{text} содержит описание
действия для диаграммы и не влияет на исполнение.

\subsection{Переходы}

Переходы объявляются с помощью \code{transition}.
Каждый переход связывает исходное состояние,
команду, действия и целевое состояние.

Например, \code{turn_on} срабатывает при получении
команды \code{switch_on} в состоянии \code{off}.
Он увеличивает счётчик включений и переводит
лампу в состояние \code{on}.

Переход \code{turn_off} возвращает лампу в состояние
\code{off}. Никаких дополнительных действий
он не выполняет, поэтому параметр \code{actions}
не указан.

\section{Сборка программы}

Объявленный класс описывает поведение лампы,
но для исполнения нужно создать автоматный объект
и включить его в программу.

\begin{lstlisting}[language=Python,numbers=none]
def build_program() -> tuple[Program, AutomatonObject]:
    lamp = AutomatonObject("lamp", Lamp, LampContext())
    return Program(objects=(lamp,)), lamp
\end{lstlisting}

\code{AutomatonObject} связывает имя объекта, класс
автомата и его начальный контекст. Благодаря этому
один класс можно использовать для создания
нескольких независимых объектов.

\code{Program} объединяет автоматные объекты
в одну систему. Пока программа содержит только
одну лампу, но позднее мы добавим взаимодействие
между несколькими автоматами.

Функцию \code{build_program()} удобно использовать
как отдельную точку сборки системы. Она позволяет
отделить создание объектов от их исполнения.

\section{Запуск программы}

Исполнением программы управляет \code{Runtime}.
Он хранит текущие состояния автоматов, их контексты
и обрабатывает поступающие команды.

Создадим исполнитель и проверим начальное состояние:

\begin{lstlisting}[language=Python,numbers=none]
program, lamp = build_program()
runtime = Runtime(program)

snapshot = runtime.snapshot(lamp)

assert snapshot.state is Lamp.off
assert snapshot.context.switch_on_count == 0
\end{lstlisting}

Метод \code{snapshot} возвращает снимок текущего
состояния объекта и его контекста. При запуске
лампа выключена, а счётчик включений равен нулю.

\subsection{Отправка команд}

Чтобы включить лампу, отправим команду
\code{switch_on} и запустим обработку:

\begin{lstlisting}[language=Python,numbers=none]
runtime.send(lamp, Lamp.switch_on)
runtime.run()

snapshot = runtime.snapshot(lamp)

assert snapshot.state is Lamp.on
assert snapshot.context.switch_on_count == 1
\end{lstlisting}

Метод \code{send} помещает команду в очередь,
указывая её получателя. Метод \code{run}
обрабатывает накопленные команды.

Runtime выбирает переход \code{turn_on},
выполняет его действие и переводит лампу
в состояние \code{on}. Счётчик включений
увеличивается до единицы.

Теперь выключим лампу:

\begin{lstlisting}[language=Python,numbers=none]
runtime.send(lamp, Lamp.switch_off)
runtime.run()

snapshot = runtime.snapshot(lamp)

assert snapshot.state is Lamp.off
assert snapshot.context.switch_on_count == 1
\end{lstlisting}

Лампа возвращается в состояние \code{off},
но счётчик не изменяется, поскольку переход
\code{turn_off} не содержит действий.

Попробуем выключить лампу ещё раз:

\begin{lstlisting}[language=Python,numbers=none]
runtime.send(lamp, Lamp.switch_off)
runtime.run()
\end{lstlisting}

Лампа уже находится в состоянии \code{off}, а переход
для команды \code{switch_off} определён только из
состояния \code{on}.

Поэтому Runtime не сможет обработать команду и вызовет
исключение \code{NoTransitionError}. По умолчанию такое
поведение считается ошибкой исполнения.

При этом состояние лампы не изменится.

\begin{importantbox}
Команды не вызываются как обычные методы Python.
Они отправляются через Runtime, который выбирает
переход с учётом текущего состояния автомата.
\end{importantbox}

\begin{trybox}
Добавьте команду \code{toggle}, которая переключает
лампу в противоположное состояние.

Подсказка: для одной команды можно объявить два
перехода с разными исходными состояниями.
\end{trybox}
